\documentclass[12pt, twoside, english]{book}
\usepackage[letterpaper, top=0.75in, bottom=0.75in, left=0.5in, right=0.5in]{geometry}
\setlength{\unitlength}{1in}
\setlength{\headheight}{15pt}
\setlength{\footskip}{36.135pt}
%\usepackage[cm]{fullpage}
\usepackage{fancyhdr}
\usepackage{tikz}
\pagestyle{fancy}
\usepackage{eso-pic}
\usepackage{verbatim}
\usepackage{afterpage}
\usepackage{babel}
\usepackage{hyperref}
\RequirePackage{lineno} 
\usepackage[absolute, overlay]{textpos}
%\usepackage{showframe}
%\usepackage[utf8]{inputenc}

%location definition
\newcommand{\currentpart}{}
\let\oldpart\part
\renewcommand{\part}[1]{\oldpart{#1}\renewcommand{\currentpart}{#1}}
\newcommand{\currentchapter}{}
\let\oldchapter\chapter
\renewcommand{\chapter}[1]{\oldchapter{#1}\renewcommand{\currentchapter}{#1}}
\newcommand{\currentsection}{}
\let\oldsection\section
\renewcommand{\section}[1]{\oldsection{#1}\renewcommand{\currentsection}{#1}}

%fonts
\renewcommand*{\familydefault}{\sfdefault}
\newfont{\titlefont}{rphvbo scaled 3000}
\newfont{\datefont}{rphvbo scaled 2500}
\newfont{\subfont}{rphvbo scaled 2000}

%colors
\definecolor{purple}{RGB}{154,1,205}
\definecolor{gray}{gray}{0.9}
\color{purple}

%definitions
\title{}
\author{}
\setcounter{chapter}{1}
\renewcommand{\thesection}{}
\renewcommand{\thesubsection}{}
\makeatletter
\def\@seccntformat#1{\rlap{\hskip-36pt\csname the#1\endcsname}}
\makeatother

%formats
\newcommand{\codefull}[1]{\\\\\indent\colorbox{gray}{\texttt{\parbox{\textwidth}{#1}}\hspace{-0.5in}}\\\\}
\newcommand{\code}[1]{\colorbox{gray}{\texttt{#1}}}
\newcommand{\sla}{$\backslash$}

%titlepage
\newcommand{\createtitle}
{
	\afterpage{
	\begin{textblock*}{8.5in}[0.5,0.5](5.9in,4in)
		\includegraphics[height=1.25in]{../../assets/LASAlogo_white}
	\end{textblock*}
	\begin{textblock*}{7.5in}[0,0](0.5in,5.5in)
		\centering
		{\color{white}
		{\titlefont
			\currentchapter\\
		}
		\vspace{10 mm}
		{\datefont
			\today\\
		}}
	\end{textblock*}
	
	\begin{textblock*}{7.5in}[0,0](0.5in,10in)
			\centering
			{\color{white}
			{\subfont
				Part of \currentpart
			}}
			
		\end{textblock*}
	\pagecolor{purple}
	\fancyfoot[]{}
	\fancyhead{}
	\color{black}
	\newpage}
	\maketitle
	%second blank page
	%\pagecolor{purple}
	%spacing
	%\newpage
	
	\setupheaders
}

%setup headers
\newcommand{\setupheaders}
{
	\fancyfoot{\color{black}}
	\fancyfoot[LO]{
		\includegraphics[height=0.5in]{../../assets/LASAlogo}
	}
	\fancyfoot[RE]{
		\today
	}
	\fancyhead{\color{black}}
	\fancyhead[RE]{
		\color{black}
		\currentchapter
	}
	\fancyhead[LO]{
		\color{black}
		\currentpart
	}
	
	\pagecolor{white}
	\color{black}
}

\begin{document}
%PART AND CHAPTER DETAILS
\renewcommand{\currentpart}{LASA Robotics Technical Documentation}
\renewcommand{\currentchapter}{Writing Technical and Notebook Documentation in LaTeX}
\createtitle{../../assets}

\section{What is LaTeX?}
\LaTeX{} (pronounced \emph{\textbf{LAY}-tek} or \emph{\textbf{LAH}-tek}) is a text formatting tool used globally for publication of scientific documents and online articles (including Wikipedia).  \LaTeX{} \emph{.tex} documents output as every format from \emph{PDF} to \emph{HTML}.

\section{Why LaTeX?}
\LaTeX{} allows for more flexible and consistent typesetting than other processors such as Word or InDesign.  Although \LaTeX{} requires users with some experience, most of \LaTeX{} can be done without much difficulty.  For LASA Robotics, most of the typesetting work has been done for you.  All you have to know are the basics.

\section{Installing LaTeX}
If \LaTeX{} is already installed on your machine, skip this step.
\subsection{Windows}
\subsubsection{1. Installing MikTeX}
Download the latest version of Mik\TeX{} from \url{http://www.miktex.org/download}.  Select 'Other Downloads' and find a \textbf{Non-basic} installation.  When downloading, select a \textbf{full installation}, which could take several hours but downloads all necessary packages.  When complete, run the wizard again to install downloaded packages.  As this take some time, it is recommended to install Mik\TeX{} on a flash drive and bring it with you.
\subsubsection{2. Installing TeXstudio}
Download and install the latest version of TeXstudio \textbf{after MikTex is downloaded and installed} from \url{http://texstudio.sourceforge.net}, which takes several minutes.  After this completes, you're good to go!

\subsection{Linux Debian}
Run \\\\\codefull{sudo apt-get install gedit-latex-plugin texlive-fonts-recommended latex-beamer\\ texpower texlive-pictures texlive-latex-extra texpower-examples imagemagick}\\ to install the GUI and all packages for \LaTeX{}.  Other distributions, including Redhat, may need alternate configurations (note that Debian-based Linux operating systems such as Ubuntu are included in this documentation.)

\subsection{Mac}
Good luck!

\newpage

\section{Types of Documentation}
There are two types of documentation written in LASA Robotics: \textbf{Technical} and \textbf{Notebook}. Although similar, there are different templates for each and both are included in the final notebook.\\\\
Technical documentation is written for \textbf{someone on the team to read}. For example, a coder would write a document describing how a piece of code works and how to use it or explain to the drivers how to set up the autonomous program \emph{in technical documentation} (this is an example of a technical doc).\\\\
Notebook documentation is written to \textbf{describe a process for a judge to read}.  For example, daily logs, email records, and goofy team photos would fit in \emph{notebook documentation}.\\\\
Follows is the tutorial for setting up either technical or notebook documentation.  If you have already completed this, scroll down and copy and paste the technical/notebook documentation examples at the bottom.

\subsection{Technical Documentation}
The first step of any documentation is to create a .tex file.  Create one using your document and save it somewhere safe.\\
Then, we begin writing.  Import the technical template like so on the first line of your .tex document: \\\\\codefull{\sla input\{../template\_technical\}}\\ The location of the file may be different, but make sure you have \code{template\_technical.tex} as the input file.

\subsection{Notebook Documentation}
As with technical documentation, the first step of any is to create a .tex file and save it somewhere safe.\\
Then, we import the notebook template like so on the first line of your .tex document:\\\\ \codefull{\sla input\{../template\_notebook\}}\\ The location of the file may be different, but make sure you have \code{template\_notebook.tex} as the input file.

\section{Setting Up the Document}
Setting up a \LaTeX{} document is easy.  Add the following line of code so your document looks like this: \\\\\codefull{\sla input\{../template\_******\}\\\sla begin\{document\}\\\\\sla end\{document\}}\\\\  Notice that all \LaTeX{} commands start with a \code{\sla}.\newpage
In the \code{\sla begin\{document\} ... \sla end\{document\}} section, add the following.
\\\\\codefull{\color{red}{\%PART AND CHAPTER DETAILS - REPLACE HIGHLIGHTED!!!}\color{black}\\\sla renewcommand\{\sla currentpart\}\{LASA Robotics {\color{red}Technical} Documentation\}\\\sla renewcommand\{\sla currentchapter\}\{{\color{red}Topic, ex.Writing Documentation in LaTeX}\}\\\sla createtitle\{{\color{red}relative location of assets folder}\}}\\\\
You should be done with initialization now.  For future reference, the default templates for technical and notebook documentation are below.

\section{Default Templates}
Feel free to copy and paste, then replace anything in red.
\subsection{Technical Documentation}
\codefull{\sla input\{{\color{red} ../template\_technical}\}\space\space\%relative path to template\_technical.tex\\\\
\sla begin\{document\}\\
\sla renewcommand\{\sla currentpart\}\{LASA Robotics Technical Documentation\}\\
\sla renewcommand\{\sla currentchapter\}\{{\color{red}Topic}\}\\
\sla createtitle\{{\color{red}../../assets}\}\space\space\%relative path to assets folder\\\\
\%-----INSERT CONTENT HERE-----\%\\\\\\
\sla end\{document\}}
\subsection{Notebook Documentation}
\codefull{\sla input\{{\color{red} ../template\_notebook}\}\space\space\%relative path to template\_notebook.tex\\\\
\sla begin\{document\}\\
\sla renewcommand\{\sla currentpart\}\{LASA Robotics Notebook Documentation\}\\
\sla renewcommand\{\sla currentchapter\}\{{\color{red}Topic}\}\\
\sla createtitle\{{\color{red}../../assets}\}\space\space\%relative path to assets folder\\\\
\%-----INSERT CONTENT HERE-----\%\\\\\\
\sla end\{document\}}\newpage

\section{Syntax Guide}
\subsection{Basic Symbols}
\code{\sla command}: Indicates a command\\
\code{\sla\sla}: Indicates a new line\\
\code{\sla\sla\sla\sla}: Indicates a new paragraph\\
\code{\sla sla}: The \code{\sla} symbol\\
\code{\%}: Indicates a comment not read by \LaTeX{}\\
\code{\%\%}: The \code{\%} symbol\\
\code{\$ ... \$}: Inline math mode (ex: $I^2 C$)

\subsection{Document Structure}
There are five primary document structures LaTeX uses in its formatting, from greatest to smallest.\\\\
\code{\sla section\{title\}}: Creates a new section\\
\code{\sla subsection\{title\}}: Creates a new subsection\\
\code{\sla subsubsection\{title\}}: Creates a new subsubsection\\
\code{\sla paragraph\{title\}}: Creates a new paragraph\\
\code{\sla subparagraph\{title\}}: Creates a new subparagraph
\subsection{Font Manipulation}
Surround text and command like this: \code{\{\sla emph Hello world!\}} to only affect small portions of text.\\
\code{\sla emph}: \emph{Emphasis}\\
\code{\sla textbf}: \textbf{Boldface}\\
\code{\sla underline}: \underline{Underline}\\
\code{\sla texttt}: \texttt{Typewriter}\\
\code{\sla code}: \code{Codebox}\\
\code{\sla url}: URL - \url{http://lasarobotics.org}\\
\code{\sla left}: Left-aligned\\
\code{\sla right}: Right-aligned\\
\code{\sla indent}: Indent text like a paragraph\\
\code{\sla tiny}: {\tiny Tiny text}\\
\code{\sla small}: {\small Small text}\\
\code{\sla normalsize}: {\normalsize Normal text}\\
\code{\sla large}, \code{\sla Large}, \code{\sla LARGE}, \code{\sla HUGE}: Even bigger fonts

\subsection{Environments}
Environments run \LaTeX{}'s powerfu
l processing.  With these, you can create anything from centered text to lists.  Remember that wherever there is a \code{\sla begin}, there must be an \code{\sla end}.
\subsubsection{Text Manipulation}
\code{\sla begin\{center\}}: Centers large amounts of text\\
\code{\sla begin\{flushright\}}: Right aligns large amounts of text\\
\code{\sla begin\{flushleft\}}: Left aligns large amounts of text
\subsubsection{Bulleted Lists}
\codefull{\sla begin\{itemize\}\\\sla item Hello\\\sla end\{itemize\}}
\subsubsection{Enumerated Lists}
\codefull{\sla begin\{enumerate\}\\\sla item Hello\\\sla end\{enumerate\}}
\subsubsection{Paragraph Codeboxes}
Paragraph codeboxes may need padding (\code{\sla\sla\sla\sla}) to appear correctly on the page.  They allow for newlines (\code{\sla\sla}) unlike inline \code{\sla code} boxes.
\\\\\codefull{\sla codefull\{Text\}}

\subsection{Graphics}
\LaTeX{} can automagically include JPG, PNG, and EPS (vector) graphics.\\
\code{\sla includegraphics\{relative location\}}\\\\
If you wish, you can even set the width and height of the graphic like this:\\
\code{\sla includegraphics[width=1in, height=350px]\{relative location\}}\\\\

\section{That's all folks!}
Once you give the basic syntax a try, you are ready to write new docs!\\Congratulations and enjoy writing beautiful documents!

\end{document}